\documentclass[11pt]{article}

\input{../shared/preamble.tex}

\usepackage{amsmath}
\usepackage{booktabs}
\usepackage{cleveref}

\title{Hierarchical Risk Parity Has No Regime of Advantage}
\author{Peter Cotton}
\date{September 21, 2026}

\newcommand{\E}{\mathbb{E}}
\newcommand{\R}{\mathbb{R}}
\newcommand{\one}{\mathbf{1}}

\newtheorem{proposition}{Proposition}
\newtheorem{lemma}{Lemma}
\newtheorem{remark}{Remark}

\begin{document}
\maketitle

\begin{abstract}
Hierarchical risk parity is defended on the regime where the sample
covariance cannot be inverted. We test that defence on its own terms, with
the population covariance known by construction so realized risk is exact,
and find no interval of sample size where the method is the right answer.
Below a tenth of an observation per asset it is beaten by inverse-variance
weighting, which inverts nothing either and consults no tree. Above parity it
is beaten by any regularised optimizer, in 78 percent of five hundred
randomly drawn dependence structures, and by long-only minimum variance on
the raw sample covariance in 76 percent, so the no-short-sale constraint
alone closes the gap. We also show that no shrinkage class small enough to be
interpretable can reproduce the method, by a counting argument confirmed at
three universe sizes, and that nested clustered optimization on the same
clusters tracks the shrinkage frontier while hierarchical risk parity does
not, which locates the cost in the inverse-variance split rather than in the
hierarchy.
\end{abstract}

\section{Introduction}

The case for hierarchical risk parity \parencite{lopezdeprado2016} rests on a comparison. The method is
set against mean-variance optimization on a raw sample covariance, it wins,
and the win is attributed to the allocation rule. Two things make that
comparison hard to read. The rule and the covariance estimate are varied
together, and the benchmark is an optimizer at its worst.

This paper fixes both and asks whether an interval of sample size survives in
which the hierarchical rule is the right choice. The population covariance is
known by construction throughout, so realized risk is $w^\top \Sigma w$
exactly rather than a backtest estimate, and every method sees the identical
sample on every draw. Where an optimizer appears it is long-only, so a
comparison against a long-only heuristic does not measure the constraint.

The answer is that no such interval exists. The window closes from both
sides, and for different reasons at each end.

\section{No parsimonious shrinkage reproduces the rule}

A natural way to read a heuristic allocator is as an optimizer fed a modified
covariance. That reading is always available and carries no information.

\begin{proposition}[Every portfolio is a minimum-variance portfolio]
\label{prop:implied}
Let $w \in \R^n$ satisfy $\one^\top w = 1$. Put $u = w/\|w\|$,
$Q = I - uu^\top$, and $\Sigma' = \one\one^\top + \varepsilon\,QMQ^\top$ for
any $\varepsilon > 0$ and symmetric positive definite $M$. Then
$\Sigma' \succ 0$ and $\Sigma' w = \one$, so $w$ is the global
minimum-variance portfolio of $\Sigma'$.
\end{proposition}

Long-short portfolios are included, and the construction asks nothing of $w$
beyond a unit budget. The question with content is therefore not whether an
implied covariance exists but whether one exists in a class small enough to
mean something.

Class size settles that before any data is seen. A family with $k$ parameters
traces at most a $k$-dimensional set of weight vectors at a fixed sample
covariance, and the weight vector has $n-1$ free coordinates, so for
$k < n-1$ the family generically misses. \Cref{tab:cliff} confirms the
threshold empirically, fitting a taper $S \circ (1 + B\theta)$ for a random
orthonormal basis $B$ of $k$ directions, with $\theta$ refitted for every
market.

\begin{table}[ht]
\centering
\begin{tabular}{rrrrr}
\toprule
assets & last $k$ that misses & its $L^1$ & first $k$ that fits & its $L^1$ \\
\midrule
16 & 15 & $1.2\times 10^{-2}$ & 16 & $3.1\times 10^{-10}$ \\
24 & 22 & $9.6\times 10^{-3}$ & 23 & $2.0\times 10^{-7}$ \\
40 & 38 & $1.1\times 10^{-2}$ & 39 & $1.9\times 10^{-4}$ \\
\bottomrule
\end{tabular}
\caption{The cliff sits at $n-1$ and tracks it, dropping four to eight orders
of magnitude in one step.}
\label{tab:cliff}
\end{table}

So the failure of any particular small family is a counting fact rather than
a failure of search. Among named families the closest is near-diagonal: a
two-parameter taper that keeps about half the correlation between assets
pairing at the finest split and none above it, leaving a fifth of the
portfolio misplaced. Shrinkage toward constant correlation, which is the best
estimator in that comparison, is the furthest from the rule.

A control confirms the miss is informative rather than generic. Fitted to a
portfolio that genuinely is minimum variance on a diagonally shrunk
covariance, the diagonal family recovers it exactly; fitted to the
hierarchical rule it misses by 0.22.

\section{The window closes from above}

\Cref{tab:robust} reports five hundred randomly drawn structures. The family,
the number and sizes of clusters, correlation levels, extra common factors,
volatility dispersion, the universe size, the sample ratio and the tail
weight are all resampled per draw. The shrinkage is a single fixed setting,
never tuned.

\begin{table}[ht]
\centering
\begin{tabular}{lrr}
\toprule
comparison & wins & median variance ratio \\
\midrule
fixed taper beats HRP & 78\% & 0.740 \\
nested clustered optimization beats HRP & 77\% & 0.738 \\
long-only minimum variance on the raw sample beats HRP & 76\% & 0.726 \\
fixed taper beats nested clustered optimization & 43\% & 1.005 \\
\bottomrule
\end{tabular}
\caption{Five hundred randomised dependence structures.}
\label{tab:robust}
\end{table}

The third row is the one to dwell on. An optimizer whose only regularisation
is the long-only constraint, which is shrinkage in the sense of
\textcite{jagannathanma2003}, beats the hierarchical rule three times in four.

The advantage holds in every structural family and is strongest where the
method's own assumption is true: 89 percent on blocks with factors, 88
percent on blocks, 83 percent on a nested hierarchy, 81 on equicorrelation,
71 on a pure factor market and 60 on an unstructured one.

\section{The window closes from below}

\Cref{tab:deep} continues into the rank-deficient corner, which is the case
the method is usually defended on.

\begin{table}[ht]
\centering
\begin{tabular}{rrrrrr}
\toprule
$T/n$ & Ledoit--Wolf & $\phi=0.5$ taper & block diagonal & NCO & inverse variance \\
\midrule
0.10 & 50\% / 1.002 & 46\% / 1.033 & 49\% / 1.006 & 36\% / 1.171 & 59\% / 0.958 \\
0.20 & 47\% / 1.027 & 45\% / 1.052 & 45\% / 1.049 & 44\% / 1.094 & 52\% / 0.996 \\
0.30 & 54\% / 0.974 & 54\% / 0.930 & 55\% / 0.954 & 55\% / 0.920 & 44\% / 1.015 \\
0.50 & 65\% / 0.879 & 63\% / 0.811 & 63\% / 0.861 & 64\% / 0.804 & 40\% / 1.018 \\
1.00 & 77\% / 0.754 & 75\% / 0.680 & 75\% / 0.687 & 75\% / 0.666 & 35\% / 1.040 \\
\bottomrule
\end{tabular}
\caption{Share of draws beating the hierarchical rule, and the median ratio
of variance to its own. Universes of forty and one hundred assets.}
\label{tab:deep}
\end{table}

The method has a genuine regime, and it is narrower than advertised. At a
tenth of an observation per asset it edges the taper by three percent and
nested clustered optimization by seventeen, because that method still has to
invert its cluster blocks and cannot.

The last column withdraws it. In the same corner, inverse-variance weighting
beats the hierarchical rule in 59 percent of draws and is four percent better
at the median. That rule inverts nothing either, consults no tree, and is
four lines of code. So at the point of maximum rank deficiency the hierarchy
earns nothing over the crudest construction that shares its one structural
advantage, and by $T/n = 0.3$ inverse variance has fallen behind while
everything else has passed.

\section{The hierarchy is not what fails}

Nested clustered optimization runs the same clustering and solves minimum
variance at both levels. On five blocks at ten observations per asset the
hierarchical rule realizes 0.1661 and it realizes 0.0892. The one-parameter
taper that dominates the first, winning every draw by a factor of about two,
beats the second in 70 percent of draws by six parts in ten thousand.

One sits on the shrinkage frontier and the other does not, from the same tree
and the same estimate. What separates them is the split rule, and the
inverse-variance split is what the hierarchical method spends its structure
on.

\section{Scope}

Two limits bound the reading above, and the second is not a footnote.

The comparisons are simulated. The population covariance has to be known for
realized risk to be exact, which is the property that makes the question
answerable at all, and it is bought by giving up real data.

More importantly, the dominance of a shrunk estimate is a property of a
stationary market. On a regime-switching market, where a crash regime fires
eight percent of the time with an independent dependence structure over a
permuted labelling, Ledoit--Wolf minimum variance is 1.21 times the
hierarchical rule's variance at a tenth coverage and 1.10 at three tenths.
Shrinking toward a single target cannot represent two dependence structures.
The claim that a simple shrinkage rule dominates holds where the covariance
is fixed and reverses where it switches.

\section{Closing}

There is no interval of sample size on which hierarchical risk parity is the
right answer. Below roughly a fifth of an observation per asset, not
inverting is what helps, and inverse-variance weighting does that better.
Above it, inverting properly is what helps, and any regularised optimizer
does that better, including one whose only regularisation is a constraint
most practitioners impose anyway. Between the two the method is beaten from
both directions at once.

The counting argument explains why the usual defences are hard to check. No
shrinkage class small enough to be interpretable can reproduce the rule, so
there is no compact estimator-side statement of what it does, and the
comparison has to be made on out-of-sample risk with the truth known. Once it
is, the hierarchy turns out not to be the problem: the same clustering with a
proper solve at both levels tracks the frontier. The cost is in the split.

\printbibliography

\end{document}
